

\begin{table*}[t!]
\centering
\small
\setlength{\tabcolsep}{4pt}
\resizebox{\linewidth}{!}{
    \begin{tabular}{ll|ccccc|ccccc}
    \toprule
    & & \multicolumn{5}{c|}{\textbf{Conflict-VQA (Acc~$\uparrow$)}} & \multicolumn{5}{c}{\textbf{PhD-icc (Acc~$\uparrow$)}} \\
    \textbf{Decoding} & \textbf{Method} & InstructBLIP & LLaVA-1.5 & LLaVA-1.6 & Qwen2.5-VL & Qwen3-VL & InstructBLIP & LLaVA-1.5 & LLaVA-1.6 & Qwen2.5-VL & Qwen3-VL \\
    \midrule
    \multirow{4}{*}{Greedy}
     & Vanilla & \num{39.61} & \num{36.20} & \num{42.35} & \num{65.61} & \num{75.58} & \num{41.08} & \num{27.72} & \num{28.23} & \num{51.30} & \num{60.32} \\
     & DoLa    & \num{32.72} & \num{40.10} & \num{44.59} & \num{66.22} & \num{76.06} & \num{32.65} & \num{29.78} & \num{28.73} & \num{51.08} & \num{58.00} \\
     & DeCo    & \num{41.31} & \num{38.03} & \num{43.38} & \num{63.79} & \num{75.33} & \num{42.12} & \num{31.00} & \num{33.65} & \num{51.65} & \num{62.25} \\
     \rowcolor{bestrow}
     &\textbf{ CALRD (Ours) }& \resg{44.12}{4.5} & \resg{41.90}{5.7} & \resg{49.51}{7.2} & \resg{70.15}{4.5} & \resg{77.42}{1.8} & \resg{49.62}{8.5} & \resg{36.35}{8.6} & \resg{36.52}{8.3} & \resg{56.20}{4.9} & \resg{63.68}{3.4} \\
    \midrule
    \multirow{4}{*}{Beam}
     & Vanilla & \num{39.98} & \num{38.88} & \num{43.99} & \num{65.49} & \num{75.94} & \num{40.73} & \num{28.25} & \num{28.80} & \num{51.11} & \num{62.53} \\
     & OPERA   & \num{40.12} & \num{39.00} & \num{43.74} & \num{68.49} & \num{76.43} & \num{39.78} & \num{28.93} & \num{29.95} & \num{53.70} & \num{62.62} \\
     & DeCo    & \num{40.83} & \num{38.15} & \num{43.01} & \num{63.55} & \num{75.46} & \num{42.15} & \num{31.07} & \num{34.63} & \num{51.13} & \num{61.25} \\
     \rowcolor{bestrow}
     & \textbf{CALRD (Ours) }& \resg{44.49}{4.5} & \resg{44.47}{5.6} & \resg{50.79}{6.8} & \resg{70.25}{4.8} & \resg{79.34}{3.4} & \resg{49.75}{9.0} & \resg{36.58}{8.3} & \resg{37.40}{8.6} & \resg{56.70}{5.6} & \resg{64.60}{2.1} \\
    \midrule
    \multirow{4}{*}{Nucleus}
     & Vanilla & \num{23.94} & \num{38.76} & \num{41.92} & \num{61.48} & \num{73.39} & \num{39.23} & \num{29.33} & \num{30.08} & \num{51.05} & \num{62.23} \\
     & VCD     & \num{23.05} & \num{42.40} & \num{46.35} & \num{64.04} & \num{75.50} & \num{40.01} & \num{31.87} & \num{31.70} & \num{53.83} & \num{63.40} \\
     & DeCo    & \num{24.91} & \num{38.45} & \num{45.44} & \num{63.55} & \num{74.79} & \num{41.83} & \num{31.25} & \num{32.90} & \num{51.62} & \num{61.50} \\
     \rowcolor{bestrow}
     &\textbf{ CALRD (Ours)} & \resg{28.07}{4.1} & \resg{43.40}{4.6} & \resg{48.72}{6.8} & \resg{68.91}{7.4} & \resg{76.55}{3.2} & \resg{48.62}{9.4} & \resg{34.95}{5.6} & \resg{34.12}{4.0} & \resg{55.83}{4.8} & \resg{64.52}{2.3} \\
    \bottomrule
    \end{tabular}
}
\caption{Performance comparison on knowledge conflict benchmarks, where textual context contradicts visual evidence. We evaluate five MLLMs with three decoding strategies. Best results are in \textbf{bold}. {\footnotesize\textcolor{red}{↑}}: improvement over Vanilla.}
\label{tab:conflict}
\end{table*}

\section{Experiments}
\label{sec:experiments}


We evaluate CALRD on two questions: (1) Does it improve conflict resolution? (2) Does it hurt non-conflict scenarios? Experiments cover five MLLMs with different architectures and capability levels. For each model, we report results on conflict benchmarks where context contradicts the image, as well as standard hallucination benchmarks.

\subsection{Experimental Setup}

\textbf{Models.} We evaluate five MLLMs: InstructBLIP-7B~\cite{instructblip2023}, LLaVA-1.5-7B, LLaVA-1.6-7B~\cite{llava-1.5}, Qwen2.5-VL-7B~\cite{bai2025qwen2}, and Qwen3-VL-8B~\cite{bai2025qwen3vltechnicalreport}. These models span different architectures and capability levels.

\textbf{Benchmarks.} We use two types of benchmarks.

For conflict resolution, we use \textbf{Conflict-VQA} (Section~\ref{sec:setup}), which contains 5,969 samples requiring open-ended answers, and \textbf{PhD-icc}~\cite{PhD_ChatGPT}, the incorrect-context subset of PhD containing 16,844 samples with yes/no question pairs. Both benchmarks feature contexts that contradict the image. We report accuracy.

For standard hallucination without explicit conflicts, we use \textbf{POPE}~\cite{pope2023} and \textbf{CHAIR}~\cite{Object_Hallucination}. POPE is a polling-based evaluation that probes object hallucination by asking ``Is there a \texttt{<object>} in the image?'' across three splits: random, popular, and adversarial. It covers 500 MSCOCO images with six questions per image for each split. We report F1 score. CHAIR measures caption hallucination by comparing mentioned objects against ground-truth annotations, reporting instance-level $\text{C}_I$ and sentence-level $\text{C}_S$ (lower is better). Following~\cite{huang2024opera}, we evaluate on 500 MSCOCO images with the prompt ``Please describe this image in detail.''


\textbf{Baselines.} We compare against four representative methods. DoLa~\cite{dola2024} contrasts logits from later layers against earlier layers to surface factual knowledge. VCD~\cite{vcd2024} contrasts outputs from original and distorted visual inputs to reduce over-reliance on language priors. OPERA~\cite{huang2024opera} penalizes over-trust patterns in self-attention and uses rollback to re-select tokens. DeCo~\cite{wang2024deco} adaptively selects preceding layers where visual information is stronger and integrates their predictions into the final output. We test CALRD with greedy, beam search, and nucleus sampling. All methods use default hyperparameters. Experiments run on NVIDIA H100 GPUs.

\subsection{Main Results}

\textbf{Does CALRD improve conflict resolution?}

Table~\ref{tab:conflict} shows results on conflict benchmarks. CALRD outperforms all baselines across models and decoders. Gains are largest on PhD-icc: LLaVA-1.5 improves from 27.72\% to 36.35\% (+8.63\%), LLaVA-1.6 from 28.23\% to 36.52\% (+8.29\%), and InstructBLIP from 41.08\% to 49.62\% (+8.54\%). On Conflict-VQA, improvements range from 4--9\%.

\textbf{Why are the gains larger on PhD-icc?} The two benchmarks differ in question format: PhD-icc uses binary yes/no questions, while Conflict-VQA requires open-ended answers. We suspect the binary format creates conditions where textual bias is more pronounced. When context strongly suggests ``yes'' or ``no,'' the model may be more easily swayed away from visual evidence. In contrast, open-ended questions require generating specific content, making it harder for context alone to override perception.

\noindent \textbf{Baseline accuracy and improvement.} An interesting pattern emerges when we rank models by their vanilla accuracy on PhD-icc: LLaVA-1.5 (27.72\%) $<$ LLaVA-1.6 (28.23\%) $<$ InstructBLIP (41.08\%) $<$ Qwen2.5-VL (51.30\%) $<$ Qwen3-VL (60.32\%). The gains from CALRD roughly follow the inverse order, with weaker models seeing larger improvements. This pattern is expected: stronger models already resolve more conflicts correctly in their late layers, leaving fewer overridden predictions for CALRD to recover. The adaptive nature of our intervention naturally accommodates this---when predictions survive to the output ($\rho$ remains high), $\lambda$ stays small and CALRD leaves the output largely unchanged. For Qwen3-VL, the +3.36\% absolute gain corresponds to an 8.5\% relative error reduction, indicating that even well-calibrated models contain recoverable failures.


\noindent \textbf{Comparison with other layer-based methods.} DeCo and DoLa also leverage information from earlier layers, but their performance is less consistent. DoLa improves some configurations but hurts others: on InstructBLIP, Conflict-VQA accuracy drops from 39.61\% to 32.72\%. DeCo shows more stable behavior but still underperforms CALRD across the board. One possible explanation is that these methods apply corrections more uniformly, without distinguishing cases where late-layer processing is beneficial. CALRD's detection signals allow it to intervene selectively, reducing the risk of disrupting predictions that were already on track.


\textbf{Does CALRD hurt non-conflict scenarios?}

A method that improves conflict resolution but degrades standard performance would have limited practical value. Tables~\ref{tab:pope_results} and~\ref{tab:chair_results} show that this is not the case for CALRD.

\noindent \textbf{POPE results (Table~\ref{tab:pope_results}).} CALRD maintains or improves F1 across all configurations. Under greedy decoding, InstructBLIP improves from 80.0 to 85.4, and Qwen3-VL from 88.5 to 91.6. POPE probes object hallucination without explicit visual-textual conflicts, so these results suggest that late-layer override may occur more broadly than just in conflict settings. The key is that CALRD does not intervene indiscriminately: when $\rho$ is high, indicating the prediction is stable, $\lambda$ stays near zero and the output remains unchanged.

\noindent \textbf{CHAIR results (Table~\ref{tab:chair_results}).} Captioning differs from VQA in that it requires extended sequences rather than short answers. Hallucinations accumulate over multiple tokens, making it a useful stress test. CALRD reduces both $\text{C}_I$ and $\text{C}_S$ in most configurations. InstructBLIP shows the largest improvement: $\text{C}_I$ drops from 23.7 to 14.6 (38\% relative reduction) and $\text{C}_S$ from 58.8 to 40.2. These results indicate that recomputing detection signals at each token position allows CALRD to adapt throughout the sequence, rather than a fixed correction.


\begin{table}[t!]
\centering
\resizebox{\linewidth}{!}{
    \setlength{\tabcolsep}{4pt}
    \begin{tabular}{ll|ccccc}
    \toprule
    & & \multicolumn{5}{c}{\textbf{POPE (F1~$\uparrow$)}} \\
    \cmidrule(l){3-7}
    \textbf{Dec.} & \textbf{Method} & \textbf{InstructBLIP} & \textbf{LLaVA-1.5} & \textbf{LLaVA-1.6} & \textbf{Qwen2.5-VL} & \textbf{Qwen3-VL} \\
    \midrule
    \multirow{4}{*}{Greedy}
     & Vanilla & \num{80.0} & \num{82.2} & \num{85.4} & \num{80.8} & \num{88.5} \\
     & DoLa    & \num{83.4} & \num{83.2} & \num{87.1} & \num{82.1} & \num{88.6} \\
     & DeCo    & \num{84.9} & \num{83.8} & \num{87.6} & \num{81.4} & \num{90.6} \\
      \rowcolor{bestrow}
     & \textbf{CALRD (Ours)} & \resup{85.4}{5.4} & \resup{84.5}{2.3} & \resup{88.2}{2.8} & \resup{82.3}{1.5} & \resup{91.6}{3.1} \\
    \midrule
    \multirow{4}{*}{Beam}
     & Vanilla & \num{84.4} & \num{84.9} & \num{87.1} & \num{80.7} & \num{88.6} \\
     & OPERA   & \num{84.8} & \num{85.4} & \num{87.5} & \num{80.8} & \num{89.3} \\
     & DeCo    & \num{84.9} & \num{86.7} & \num{87.9} & \num{81.4} & \num{91.2} \\
      \rowcolor{bestrow}
     & \textbf{CALRD (Ours)} & \resup{84.7}{0.3} & \resup{86.9}{2.0} & \resup{88.3}{1.2} & \resup{82.6}{1.9} & \resup{91.9}{3.3} \\
    \midrule
    \multirow{4}{*}{Nucleus}
     & Vanilla & \num{79.8} & \num{83.1} & \num{85.5} & \num{80.8} & \num{88.7} \\
     & VCD     & \num{79.9} & \num{83.1} & \num{85.6} & \num{80.9} & \num{88.9} \\
     & DeCo    & \num{81.8} & \num{84.8} & \num{87.3} & \num{81.3} & \num{91.2} \\
      \rowcolor{bestrow}
     & \textbf{CALRD (Ours)} & \resup{83.4}{3.6} & \resup{85.9}{2.8} & \resup{87.5}{2.0} & \resup{81.9}{1.1} & \resup{91.4}{2.7} \\
    \bottomrule
    \end{tabular}
}
\caption{Results on POPE object hallucination benchmark. The best results are in bold. {\footnotesize\textcolor{red}{↑}}: improvement over Vanilla.}
\label{tab:pope_results}
\end{table}

\vspace{-10pt}

\begin{table*}[t]
\centering
\small
\setlength{\tabcolsep}{3pt}
\resizebox{\linewidth}{!}{
    \begin{tabular}{ll|cc|cc|cc|cc|cc}
    \toprule
    & & \multicolumn{2}{c|}{\textbf{InstructBLIP}} & \multicolumn{2}{c|}{\textbf{LLaVA-1.5}} & \multicolumn{2}{c|}{\textbf{LLaVA-1.6}} & \multicolumn{2}{c|}{\textbf{Qwen2.5-VL}} & \multicolumn{2}{c}{\textbf{Qwen3-VL}} \\

    \textbf{Decoding} & \textbf{Method} & $\text{C}_S \downarrow$ & $\text{C}_I \downarrow$ & $\text{C}_S \downarrow$ & $\text{C}_I \downarrow$ & $\text{C}_S \downarrow$ & $\text{C}_I \downarrow$ & $\text{C}_S \downarrow$ & $\text{C}_I \downarrow$ & $\text{C}_S \downarrow$ & $\text{C}_I \downarrow$ \\

    \midrule

    \multirow{4}{*}{Greedy}

     & Vanilla & \numb{58.8} & \num{23.7} & \num{45.0} & \num{14.7} & \num{37.6} & \num{12.7} & \num{40.8} & \num{9.7} & \num{52.4} & \num{10.1} \\

     & DoLa    & \numb{48.4} & \num{15.9} & \num{47.8} & \num{13.8} & \num{35.3} & \num{8.7} & \num{36.5} & \num{9.2} & \num{55.6} & \num{9.8} \\

     & DeCo    & \numb{41.2} & \num{14.4} & \num{37.8} & \num{11.1} & \num{35.8} & \num{10.4} & \num{37.6} & \textbf{\num{9.1}} & \num{51.2} & \num{9.6} \\
 \rowcolor{bestrow}
     & \textbf{CALRD (Ours)} & \resdown{40.2}{18.6} & \resdown{14.6}{9.1} & \resdown{35.7}{9.3} & \resdown{10.3}{4.4} & \resdown{33.6}{4.0} & \resdown{9.3}{3.4} & \resdown{35.5}{5.3} & \resdown{9.5}{0.2} & \resdown{50.7}{1.7} & \resdown{9.3}{0.8} \\

    \midrule

    \multirow{4}{*}{Beam}

     & Vanilla & \numb{55.6} & \num{15.8} & \numb{48.8} & \num{13.9} & \num{36.0} & \num{12.1} & \num{38.9} & \num{9.4} & \numb{53.6} & \num{9.7} \\

     & OPERA   & \numb{46.4} & \num{14.2} & \numb{44.6} & \num{12.8} & \num{35.3} & \num{12.9} & \num{39.1} & \num{10.2} & \numb{54.7} & \num{11.3} \\

     & DeCo    & \numb{43.8} & \num{12.7} & \numb{32.0} & \num{9.7} & \num{34.4} & \num{9.0} & \num{38.2} & \num{8.5} & \numb{52.0} & \num{9.3} \\
 \rowcolor{bestrow}
     & \textbf{CALRD (Ours)} & \resdown{38.4}{17.2} & \resdown{11.4}{4.4} & \resdown{34.7}{14.1} & \resdown{10.9}{3.0} & \resdown{32.8}{3.2} & \resdown{8.7}{3.4} & \resdown{33.4}{5.5} & \resdown{8.2}{1.2} & \resdown{40.6}{13.0} & \resdown{7.9}{1.8} \\

    \midrule

    \multirow{4}{*}{Nucleus}

     & Vanilla & \numb{54.6} & \numb{24.8} & \numb{48.8} & \num{14.2} & \num{36.8} & \num{10.5} & \num{40.4} & \num{10.7} & \num{54.2} & \num{10.2} \\

     & VCD     & \numb{58.0} & \numb{17.0} & \numb{54.0} & \num{16.0} & \num{41.2} & \num{10.2} & \num{42.9} & \num{12.3} & \num{53.4} & \num{11.0} \\

     & DeCo    & \numb{43.6} & \numb{12.9} & \numb{42.8} & \num{13.2} & \num{36.1} & \num{10.1} & \num{40.4} & \num{9.5} & \num{54.8} & \num{9.9} \\
 \rowcolor{bestrow}
     & \textbf{CALRD (Ours)} & \resdown{42.7}{11.9} & \resdown{12.1}{12.7} & \resdown{37.6}{11.2} & \resdown{12.8}{1.4} & \resdown{34.8}{2.0} & \resdown{11.5}{-1.0} & \resdown{38.6}{1.8} & \resdown{9.1}{1.6} & \resdown{53.6}{0.6} & \resdown{9.5}{0.7} \\

    \bottomrule

    \end{tabular}

}

\caption{Results on CHAIR caption hallucination benchmark. Lower is better. {\footnotesize\textcolor{blue}{↓}}: reduction over Vanilla.}
\label{tab:chair_results}
\end{table*}

\begin{table}[t]
\centering
\begin{tabular}{l|cc}
\toprule
Configuration & C-VQA $\uparrow$ & POPE$\uparrow$  \\
\midrule
Vanilla (Greedy) & 42.3 & 85.4  \\
\midrule
CALRD (Full)     & \textbf{49.5} & \textbf{88.2}  \\
\quad w/o $D_{\text{conf}}$ & 47.3 & 87.5  \\
\quad w/o $\rho$  & 45.1 & 86.6  \\
\quad w/ Fixed Layer (75\%) & 45.9 & 87.1  \\
\bottomrule
\end{tabular}
\caption{Ablation study on LLaVA-1.6 with greedy decoding. We evaluate the contribution of each component.}
\label{tab:ablation}
\end{table}


\begin{figure}[t]
  \centering
  \includegraphics[width=\linewidth]{figures/efficiency.pdf}
  \caption{\textbf{Inference Efficiency on LLaVA-1.5-7B.} Latency vs. GPU memory trade-off measured on a single H100 GPU. Our CALRD achieves a superior balance, staying closest to the Vanilla baseline compared to DECO, VCD, and OPERA.}
  \label{fig:efficiency}
\end{figure}



\subsection{Analysis}

\textbf{Ablation study.} Table~\ref{tab:ablation} examines the contribution of each component using LLaVA-1.6 with greedy decoding. Removing $D_{\text{conf}}$ drops Conflict-VQA accuracy from 49.5\% to 47.3\%, while removing $\rho$ causes a larger drop to 45.1\%. This suggests that retention is the more informative signal for detecting harmful override, which makes sense: a prediction can be confident at the transition layer for various reasons, but a sharp drop in retention specifically indicates that something changed in late-layer processing.

Using a fixed layer at 75\% depth instead of dynamic $L_{\text{trans}}$ detection reduces accuracy to 45.9\%. The value corresponds to the median transition layer observed in our analysis (Section~\ref{sec:discovery2}), so this comparison uses the most representative fixed position rather than an arbitrary choice. The accuracy drop indicates that transition points vary substantially across samples, and a single fixed layer cannot capture this variation. We also observe that POPE follows a similar pattern, with both signals contributing to the overall improvement.

These results support the multiplicative design $\lambda = D_{\text{conf}} \cdot (1 - \rho)$: both signals provide useful information, but neither is sufficient alone. The product ensures that substantial intervention occurs only when a confident prediction exists and is being suppressed by late-layer processing.


\textbf{Efficiency.} Figure~\ref{fig:efficiency} compares latency and memory usage across methods using LLaVA-1.5 on a single H100 GPU. CALRD adds modest overhead: 19.8 ms/token compared to 18.3 ms/token for vanilla decoding, with only ~80 MB additional memory. VCD nearly doubles latency due to its dual-stream requirement, and OPERA is approximately \textbf{10$\times$} slower because of iterative rollback. DeCo falls in between. These results suggest CALRD is practical for deployment scenarios where efficiency matters.
